\documentclass[12pt,a4paper]{article}

\usepackage[utf8]{inputenc}
\usepackage[spanish]{babel}
\usepackage{amsmath, amssymb, amsthm}
\usepackage{graphicx}
\usepackage[margin=1in]{geometry}
\usepackage{lipsum}
\usepackage{float}
\usepackage{minted}
\usepackage{subcaption}
\usepackage{rotating}

\title{\textbf{Money Laundering Analysis}}
\date{\today}

\begin{document}

\begin{titlepage}
  \centering
  \vspace*{5cm}

  
\includegraphics[width=0.4\textwidth]{img/logo-fiuba.png}\\[1cm]

  {\Huge \textbf{Money Laundering Analysis}}\\[0.2cm]

  {
    \large \textbf{Sistemas Distribuidos}\\[0.2cm]
    \begin{tabular}{l l}
      108397 & Ordoñez Alejo\\
      107863 & Minervino Lorenzo\\
      106010 & Valsagna Federico\\
    \end{tabular}
  }

  \vfill

  {\large 1c2026}
\end{titlepage}

\tableofcontents

\newpage
\section{Arquitectura del sistema}
El sistema recibe un conjunto de datos de transacciones bancarias
desde uno o más clientes simultáneamente y devuelve los resultados de
cinco análisis distintos. La diversidad de los requerimientos, que
abarcan desde un filtro simple hasta la detección de patrones en
grafos y la consulta a servicios externos, motivó el diseño de una
arquitectura de procesamiento distribuido en pipeline, donde cada
caso de uso recorre una cadena de nodos especializados de forma
concurrente con los demás.

\section{Casos de uso}
Una petición del cliente hace que el sistema procese los 5 casos de uso:

\begin{itemize}
  \item \textbf{UC1} Cuenta de origen, cuenta de destino y monto para
    transacciones USD con monto menor a \$50.
  \item \textbf{UC2} Nombre de banco, cuenta de origen y monto de la
    máxima transacción en USD para cada banco.
  \item \textbf{UC3} Cuenta de origen y monto de transacciones USD en
    el período 2022-09-06 al
    2022-09-15 (período B), cuyo monto sea menor a un centésimo del
    promedio de monto para su formato en el período 2022-09-01 al
    2022-09-05 (período A).
  \item \textbf{UC4} Pares de cuentas que cumplan con el patrón
    \textit{scatter-gather}
    con una cuenta
    de separación y una cantidad mínima de cuentas intermedias igual
    a 5; en el período A.
  \item \textbf{UC5} Cantidad de transacciones con formato de pago
    \textit{Wire} o \textit{ACH} en
    el período A, cuyo monto en USD sea menor a 1.
\end{itemize}

\begin{figure}[h!]
  \centering
  \includegraphics[width=0.6\textwidth]{../diagrams/v2/08use-cases-diagram.png}\\[1cm]
  \caption{Diagrama de casos de uso}
\end{figure}

Los cinco casos de uso presentan niveles crecientes de complejidad:
desde un filtro directo (UC1) hasta la detección de estructuras en un
grafo de transacciones (UC4) o la normalización de moneda contra un
servicio externo (UC5). Esta escalera de complejidad da forma a la
arquitectura del sistema y se verá reflejada en cada una de las
vistas que siguen.

\newpage
\section{DAG de procesamiento}

El conjunto de pipelines puede modelarse como un DAG (\textit{Directed
Acyclic Graph}). La ausencia de ciclos garantiza que el procesamiento
siempre converge y que las dependencias entre etapas son explícitas:
un nodo solo puede procesar datos cuando sus predecesores ya los produjeron.

\subsection{\textbf{UC1} — Filtro directo}
UC1 establece el patrón base del sistema. El \textbf{Gateway} entrega el
registro completo \texttt{(origin, destination, amount, format)} y el
\textbf{Filter} aplica dos condiciones simultáneas: moneda USD y monto
menor a \$50. Los registros que las satisfacen salen como \texttt{(origin,
destination, amount)} hacia el \textbf{Join}. El campo \texttt{format} se
descarta en esta etapa porque no aporta información al resultado.

\begin{figure}[H]
  \centering
  \includegraphics[width=\textwidth,height=0.85\textheight,keepaspectratio]{../diagrams/v2/11dag-diagram-uc1.png}
  \caption{Diagrama DAG UC1}
\end{figure}

Cada transacción se evalúa de forma completamente independiente de
las demás. No hay computación entre filas, por lo que el pipeline
puede procesar los registros en cualquier orden y en paralelo sin
afectar la correctitud del resultado.

\subsection{\textbf{UC2} — Máximo por banco}
UC2 introduce dos conceptos nuevos respecto de UC1: múltiples fuentes
de datos y agregación en dos etapas. Primero se distribuyen dos
streams simultáneamente: registros de transacciones \texttt{(bank\_id,
origin, amount)} hacia el pipeline de filtrado, y registros de
cuentas \texttt{(bank\_id, bank\_name)} hacia un \textbf{GroupBy} de
resolución de
nombres que elimina registros que duplican la información con el
nombre del banco.

\begin{figure}[H]
  \centering
  \includegraphics[width=\textwidth,height=0.85\textheight,keepaspectratio]{../diagrams/v2/11dag-diagram-uc2.png}
  \caption{Diagrama DAG UC2}
\end{figure}

El pipeline de transacciones filtra por USD y luego el GroupBy
calcula el máximo de \texttt{amount} por \texttt{bank\_id} dentro de
cada partición.
El \textbf{Aggregate} consolida esos máximos parciales en el máximo global
definitivo por banco.

\subsection{\textbf{UC3} — Comparación entre períodos}
UC3 requiere comparar transacciones de dos períodos distintos: el
umbral \texttt{(amount < promedio\_de\_formato / 100)} no es un valor fijo
sino que depende del promedio global de cada formato en el período A,
que solo se conoce una vez que se procesaron todos los datos de ese período.

\begin{figure}[H]
  \centering
  \includegraphics[width=\textwidth,height=0.85\textheight,keepaspectratio]{../diagrams/v2/11dag-diagram-uc3.png}
  \caption{Diagrama DAG UC3}
\end{figure}

Primero se divide el stream de datos en dos bifurcaciones a ser
procesadas en paralelo: una con currency USD y período A y la otra
con currency USD y período B.
El pipeline del período A materializa los promedios por formato: el
GroupBy format Sum amount Count acumula suma y conteo por formato en
cada partición, y el Aggregate computa el promedio real
(\texttt{sum/count}). La salida es \texttt{(format, average)}.
El pipeline del período B produce el stream crudo de transacciones
candidatas: \texttt{(origin, amount, format)}.
El Merge combina ambas salidas usando el formato como clave, de modo
que cada transacción del período B queda asociada con el promedio de
su formato en el período A.
Un Filter final aplica la condición \texttt{amount < average/100} y produce
\texttt{(origin, amount)} para el Join.

\subsection{\textbf{UC4} — Patrón scatter-gather}
UC4 requiere detectar el patrón \textit{scatter-gather} en el grafo de
transacciones: una cuenta origen dispersa fondos hacia múltiples
intermediarios, que los concentran en una cuenta destino. Esta
detección no puede hacerse registro a registro.

\begin{figure}[H]
  \centering
  \includegraphics[width=\textwidth,height=0.85\textheight,keepaspectratio]{../diagrams/v2/11dag-diagram-uc4.png}
  \caption{Diagrama DAG UC4}
\end{figure}

Después del filtro inicial (currency USD y período A), el stream se
bifurca en dos sub-pipelines que construyen, cada uno, la misma
representación del grafo de transacciones: para cada nodo, el conjunto
de predecesores y sucesores observados \texttt{(node, predecessors,
successors)}.

La primera bifurcación calcula el \textit{grado} de cada nodo y selecciona
los nodos de \textbf{alto grado}: aquellos con al menos 5 predecesores y al
menos 5 sucesores. Un nodo que no alcanza ese grado no puede participar
de un patrón scatter-gather con 5 intermediarios, de modo que esta
etapa actúa como una poda temprana que reduce el grafo antes del conteo
costoso. La segunda bifurcación materializa el grafo completo.
El \textbf{Merge} (\textit{prune}) cruza ambas salidas y conserva únicamente
las aristas cuyos dos extremos son nodos de alto grado.

Sobre ese grafo podado, el conteo de caminos cuenta, para cada par
\texttt{(origin, destination)}, la cantidad de cuentas que sirven como
intermediario directo entre ese origen y ese destino —es decir, el
número de cuentas intermediarias del patrón scatter-gather. El Aggregate
final consolida esos conteos parciales y conserva solo los pares con un
conteo mayor o igual a 5.
El resultado \texttt{(origin, destination)} llega al Join.

\subsection{\textbf{UC5} — Conteo con conversión de moneda}
UC5 cuenta cuántas transacciones de formato Wire o ACH en el período
A tienen un monto menor a \$1 USD. Los montos pueden estar expresados
en cualquier moneda, por lo que la comparación con el umbral solo es
válida después de convertir.

\begin{figure}[H]
  \centering
  \includegraphics[width=\textwidth,height=0.85\textheight,keepaspectratio]{../diagrams/v2/11dag-diagram-uc5.png}
  \caption{Diagrama DAG UC5}
\end{figure}

El Filter inicial aplica formato (Wire o ACH) y período A para
reducir el dataset antes de la etapa de conversión.
El paso de conversión traduce cada monto a su equivalente en dólares.
El Filter siguiente aplica el umbral de \$1. Las transacciones que lo
superan llegan al GroupBy * Count, donde el \texttt{*} indica que no hay
clave de agrupación: todos los registros se colapsan en un único
contador global.
Ese conteo es el resultado final de UC5.

\section{Vista física}
\subsection{Robustez}
El diagrama de robustez muestra cómo se ejecutan los controladores en
los pipelines: qué componentes concretos existen, qué colas los
conectan, cuántas instancias corren en paralelo y qué estrategia de
ruteo usa cada etapa.

Todo el ruteo entre etapas se resuelve por \textit{afinidad}: cada enlace
es un \textit{exchange} directo de RabbitMQ con una cola por instancia
downstream, ligada por una \textit{routing key} igual a su índice de
\textit{shard}. El productor elige el shard de cada registro y publica con
esa clave, de modo que cada instancia consumidora es dueña exclusiva de
su partición. No existen \textit{working queues} con consumidores
competitivos: cada cola tiene exactamente un consumidor. Las únicas colas
de trabajo simples son las de entrada del Join y la de Responses, pero
también con un único consumidor cada una. Esta decisión es deliberada: el
modelo \textit{multi-computing} del enunciado se aprovecha repartiendo
particiones disjuntas entre instancias, no haciéndolas competir por una
misma cola.

La clave de \textit{shard} depende de la etapa. Donde el ruteo no necesita
agrupar por un atributo del dominio, el productor reparte por la
identidad del mensaje (un \textit{hash} estable derivado del identificador
del productor más su número de secuencia); donde sí hace falta que todos
los registros de un grupo caigan juntos, el shard se calcula sobre la
clave de dominio correspondiente (\texttt{bank\_id}, \texttt{format}, nodo del
grafo, etc.). Por esto último, todas las instancias de \texttt{group\_by} y
\texttt{aggregate} fijan \texttt{PYTHONHASHSEED} para que la asignación de
shards sea determinística y estable entre reinicios.

\subsubsection{UC1}
El \texttt{default\_filter} corre con múltiples instancias en paralelo
(dos en el despliegue de referencia), cada una procesando una partición
del stream de transacciones que el Gateway reparte por \textit{round-robin}.
Cada instancia aplica las dos condiciones (USD y monto menor a \$50) y,
como UC1 no tiene etapas intermedias, deposita los registros que las
satisfacen directamente en la cola de entrada del Join.

El Join no se escala por datos: está \textbf{particionado por caso de uso}.
Cada instancia de Join es dueña de las colas de uno o más casos de uso y
ninguna otra las consume, por lo que cada cola de UC tiene un único
consumidor sin necesidad de shardear por \texttt{client\_id}. Como un mismo
cliente recibe exactamente una respuesta por caso de uso, alcanza con que
la instancia dueña de ese UC reúna todos los resultados parciales del
cliente antes de emitir la respuesta.

\begin{figure}[H]
  \centering
  \includegraphics[width=0.62\textwidth]{../diagrams/v2/12robustness-diagram-uc1.png}
  \caption{Diagrama de Robustez UC1}
\end{figure}

La cola de Responses desacopla el Join del Gateway: el Join deposita
los resultados en la cola y el Gateway los consume para cerrar la
conexión correspondiente.

\subsubsection{UC2}
El Gateway publica en dos exchanges distintos: \texttt{Transactions} para
el stream de transacciones y \texttt{Accounts} para el dataset de cuentas.
Los dos pipelines —el que resuelve el nombre de banco y el que obtiene
su monto máximo— corren en paralelo. En el de transacciones, el
\texttt{default\_filter} filtra por USD, el \texttt{GroupBy bank\_id Max amount}
calcula el máximo parcial por \texttt{bank\_id} en cada partición (con shardeo
por \texttt{bank\_id} entre GroupBy y Aggregate) y el \texttt{Aggregate max
amounts} consolida el máximo global por banco. En el de cuentas, el
\texttt{GroupBy bank\_id On bank\_name} extrae las asociaciones
\texttt{(bank\_id, bank\_name)} y el \texttt{Aggregate bank names} las consolida.

El \texttt{Merge On bank\_id} cruza ambos pipelines. Corre con varias
instancias (dos en el despliegue de referencia) coordinadas en anillo. El
Aggregate de máximos hace \textit{broadcast} de su salida —cada instancia del
Merge recibe los máximos de todos los bancos—, mientras que el Aggregate
de nombres shardea por \texttt{bank\_id}, de modo que cada instancia del
Merge es dueña de un subconjunto de bancos. Así, cada banco es emitido
por exactamente una instancia del Merge: la que tiene su nombre en su
shard. El resultado se deposita en la cola de entrada del Join.

\begin{figure}[H]
  \centering
  \includegraphics[width=\textwidth,height=0.85\textheight,keepaspectratio]{../diagrams/v2/12robustness-diagram-uc2.png}
  \caption{Diagrama de Robustez UC2}
\end{figure}

\subsubsection{UC3}
El \texttt{default\_filter} produce dos exchanges de salida distintos:
\texttt{Filtered by USD and period A} y \texttt{Filtered by USD and period
B}. Además de filtrar por currency, rutea cada transacción al pipeline
correspondiente según su fecha. Una transacción no puede pertenecer a
ambos períodos, por lo que el ruteo es mutuamente exclusivo.
El pipeline del período A usa shardeo por \texttt{format} entre el GroupBy y
el Aggregate para que todos los registros del mismo formato lleguen a
la misma instancia.
Los promedios resultantes se distribuyen al Merge mediante
\textit{broadcast}: cada resultado del Aggregate se replica a todas las
instancias del Merge, mientras que las transacciones del período B se
reparten por shardeo entre esas mismas instancias. Cualquier instancia
del Merge puede recibir transacciones del período B de cualquier formato
y necesita los promedios de todos los formatos para evaluar la condición;
el shardeo enviaría cada promedio a una sola instancia, mientras que el
broadcast garantiza que todas tengan el contexto completo.
El broadcast es aceptable dado que por la naturaleza de los datos del
problema no puede haber demasiados formatos de pago; en particular,
hay 7 formatos en el dataset de ejemplo entero.

Después del Merge, un segundo Filter aplica la condición \texttt{amount <
average/100} y produce la cola final hacia el Join.

\begin{figure}[H]
  \centering
  \includegraphics[width=\textwidth,height=0.85\textheight,keepaspectratio]{../diagrams/v2/12robustness-diagram-uc3.png}
  \caption{Diagrama de Robustez UC3}
\end{figure}

\subsubsection{UC4}
Tras el \texttt{default\_filter} (USD y período A), el stream se duplica
hacia dos sub-pipelines que construyen la misma representación del grafo,
ambos con el \texttt{GroupBy Compute graph} ruteando por \texttt{node} hacia su
Aggregate. El primero, \texttt{Aggregate degree}, selecciona los nodos de
alto grado (al menos 5 predecesores y 5 sucesores) y los hace
\textit{broadcast} hacia todas las instancias del Merge. El segundo,
\texttt{Aggregate graphs}, materializa el grafo completo y lo shardea por
\texttt{node} entre las instancias del Merge.

El \texttt{Merge prune} corre con varias instancias en anillo: cada una
tiene la lista completa de nodos de alto grado (por broadcast) y su shard
del grafo (por sharding), y conserva solo las aristas cuyos dos extremos
son de alto grado. Sobre ese grafo podado, el \texttt{Count paths} cuenta
los intermediarios por par \texttt{(origin, destination)} —shardeando por
camino para repartir el conteo entre instancias— y el \texttt{Aggregate
paths} consolida los conteos parciales, conservando los pares con conteo
mayor o igual a 5 y depositándolos en la cola de entrada del Join. Por el
volumen del grafo, el conteo de caminos vuelca su estado parcial a disco,
particionado en archivos por shard.

\begin{figure}[H]
  \centering
  \includegraphics[width=\textwidth,height=0.85\textheight,keepaspectratio]{../diagrams/v2/12robustness-diagram-uc4.png}
  \caption{Diagrama de Robustez UC4}
\end{figure}

\subsubsection{UC5}
El controlador \texttt{Convert amount into USD} puede escalar
horizontalmente de forma independiente del Filter. Se optó por
permitir este escalamiento para más libertad en el despliegue del
sistema, pero la operación que se realiza en esta instancia es
puramente I/O intensive, y por tanto no requiere mayor poder de
cómputo. Se emite una sola request al servicio externo de conversión
por batch, amortizando la latencia de red sobre el conjunto del lote.
Tras la conversión, un segundo Filter aplica el umbral de \$1 USD y el
\texttt{GroupBy * Count} produce un conteo parcial por instancia. Esos
conteos parciales se depositan en la cola de entrada del Join, que los
suma en el conteo global definitivo del caso de uso.

\begin{figure}[H]
  \centering
  \includegraphics[width=\textwidth,height=0.85\textheight,keepaspectratio]{../diagrams/v2/12robustness-diagram-uc5.png}
  \caption{Diagrama de Robustez UC5}
\end{figure}

\subsubsection{Overview del sistema completo}

\begin{figure}[h!]
  \centering
  \includegraphics[width=\textwidth,height=0.9\textheight,keepaspectratio]{../diagrams/v2/12robustness-diagram.png}
  \caption{Diagrama de Robustez del sistema completo}
\end{figure}

El Filter es el nodo con una mayor cantidad de conexiones: recibe el
stream completo de transacciones y produce múltiples colas
clasificadas por período, formato y moneda. Su escalado horizontal
genera una topología \textit{many-to-many} con los nodos downstream que debe
gestionarse a nivel de exchanges y colas del broker.
Esto se modeló así porque las alternativas que surgieron son:

\begin{itemize}
  \item Segregar el filtrado en distintas instancias de filter y
    encadenarlas (por ejemplo filtrar primero por currency USD y después
    por amount menor a \$50 en el primer caso de uso).
  \item Segregar \textit{los filtros} en distintas instancias según el
    filtro a aplicar.
\end{itemize}

El problema con estas alternativas es que por la naturaleza de los
casos de uso, las cantidad de conexiones sería igual de grande, dado
que la mayoría requieren primero filtrar por currency USD. Por otro
lado, separar el filtrado en múltiples instancias hace que sea
necesario recorrer los batches varias veces en vez de sólo una y
rutear los registros en los batches de manera acorde. Además, la
segunda alternativa es equivalente en robustez a escalar
horizontalmente el único controlador Filter por el se terminó optando.

Los ruteos por clave que se realizan hacia algunos de los
controladores aggregate o Merge son necesarios para garantizar que
todos los registros de un mismo grupo se aggreguen correctamente.
Además se realizan en instancias del pipeline en las que se asegura
que la cantidad de variantes por las que se hace es tan grande como
para seguir aprovechando la estrategia \textit{multi-computing}.

El último controlador Join está \textbf{particionado por caso de uso}: cada
instancia es dueña de las colas de uno o más casos de uso (en el
despliegue de referencia, una instancia atiende UC1, otra UC3 y otra
UC2/UC4/UC5). Como los resultados parciales de un caso de uso ya
convergen todos en su única cola, no hace falta shardear por
\texttt{client\_id}: la instancia dueña del UC reúne los resultados de cada
cliente y construye la respuesta final. Esta partición permite balancear
el trabajo del Join entre instancias según el peso relativo de cada caso
de uso, en lugar de hacerlo según la cantidad de clientes.

Los nodos Merge (UC2, UC3 y UC4) son los únicos puntos donde se cruzan
pipelines que provienen de fuentes distintas. Fuera de esos casos, cada
pipeline corre de extremo a extremo sin interacción con las demás.

\newpage
\subsection{Despliegue}
Para el despliegue del sistema se intentó desacoplar el mismo por
completo de la \textit{vista lógica}: el hecho de que dos controladores
realicen un filtrado, no debería impactar en la decisión de en qué
nodos desplegarlos.
Dicho esto, en el siguiente diagrama, se muestra cómo podrían
desplegarse los procesos del pipeline de la manera más distribuida posible.

Cada controlador corre en su propio contenedor y se conecta al
\textit{broker} RabbitMQ, que actúa como middleware de comunicación por
grupos. Además de los nodos del pipeline, el despliegue incluye tres
procesos de infraestructura: el contenedor de RabbitMQ, el
\textbf{Supervisor} —que recibe los \textit{heartbeats} de todos los nodos y
revive a los que caen— y, opcionalmente, el \textbf{Chaos Monkey} usado
para inyectar fallas durante las pruebas. Cada nodo con estado monta
además un volumen propio donde persiste su \textit{checkpoint}, de modo que
al reiniciarse recupera su estado desde el mismo volumen. Estos
componentes de tolerancia a fallos se detallan en la sección
correspondiente.

\begin{sidewaysfigure}
  \centering
  \includegraphics[width=0.95\textheight,height=0.92\textwidth,keepaspectratio]{../diagrams/v2/14deployment-diagram.png}
  \caption{Diagrama de despliegue}
\end{sidewaysfigure}

\section{Vista de desarrollo}
\subsection{Paquetes}
Para la distribución de archivos fuente, se optó por mantener un
directorio por cada \textit{tipo} de controlador bajo \texttt{src/}
(\texttt{filter/}, \texttt{group\_by/}, \texttt{aggregate/}, \texttt{merge/},
\texttt{join/}, \texttt{converter/}, \texttt{gateway/} y \texttt{client/}), con su
archivo \texttt{main.py} que levanta el proceso, su correspondiente
Dockerfile y sus dependencias específicas.
Las dependencias comunes se centralizan en \texttt{common/}, que cada
Dockerfile copia junto al código del controlador. Sus subpaquetes
principales son \texttt{comms/} —la comunicación, que agrupa el middleware
sobre RabbitMQ (\texttt{middleware/}), el modelo de mensajes
(\texttt{messages/}) y la maquinaria de fin de stream (\texttt{eof\_handler/})—,
\texttt{data/} —el modelo de dominio (\texttt{Transaction}, \texttt{Account})—,
\texttt{checkpoint/} —persistencia y deduplicación— y \texttt{heartbeat/} —el
cliente de \textit{heartbeat} hacia el Supervisor—.

\begin{figure}[H]
  \centering
  \includegraphics[width=0.92\textwidth,height=0.6\textheight,keepaspectratio]{../diagrams/v2/15packages-diagram.png}
  \caption{Diagrama de paquetes}
\end{figure}

\begin{figure}[H]
  \centering
  \includegraphics[width=0.6\textwidth]{../diagrams/v2/15packages-diagram-filter.png}
  \caption{Diagrama de paquetes — \texttt{filter/}}
\end{figure}

La razón por la cuál se pueden generalizar los controladores es el
patrón de implementación por el que se optó. Cada tipo de controlador
tiene un nodo genérico —que hereda de \texttt{RingNode} (o de
\texttt{StatelessRingNode} cuando no acumula estado) y resuelve el consumo,
el ruteo por afinidad, la sincronización de fin de stream y el
\textit{checkpointing}— al que se le inyecta una \textit{función de estrategia}
que encapsula la lógica propia del caso de uso. La estrategia concreta se
elige al levantar el proceso mediante la variable de entorno
\texttt{STRATEGY}.

Por ejemplo, el código fuente para un Filter consta del nodo genérico
\begin{minted}{python}
class RingFilter(StatelessRingNode):
    def __init__(self, fn: FilterFn, ...):
        # ...
\end{minted}

de la interfaz para la función de filtro
\begin{minted}{python}
class FilterFn:
    @abstractmethod
    def filter(self, el: Message) -> Message:
        ...
\end{minted}

y de la implementación particular del filtro que se quiera correr
\begin{minted}{python}
class UC1Filter(FilterFn):
    def filter(self, el: Message) -> Message:
        # USD y monto < 50
        ...
\end{minted}

El mismo esquema se repite para los demás controladores, cada uno con su
propia interfaz de estrategia: \texttt{GroupByFn}, \texttt{AggregateFn},
\texttt{MergeFn}, \texttt{ConverterFn} y \texttt{JoinFn}. Todas operan sobre la
abstracción \texttt{Message}, de la que \texttt{Transaction} y \texttt{Account}
son los datos de dominio que viaja en los mensajes de entrada.

\newpage
\section{Vista de procesos}
\subsection{Actividades}
\subsubsection{UC1}
En este diagrama se muestra, además del pipeline necesario para
computar los resultados a retornarle al cliente, cómo se manejan los
datos que el cliente envía por batches. Esta explicación se abstrae
del resto de casos de uso, puesto a que de por sí ya presentan una
complejidad bastante más grande que la de éste.

\begin{figure}[H]
  \centering
  \includegraphics[width=\textwidth,height=0.82\textheight,keepaspectratio]{../diagrams/v2/16activity-diagram-uc1.png}
  \caption{Diagrama de actividades — UC1}
\end{figure}

\subsubsection{UC2}
\begin{sidewaysfigure}
  \centering
  \includegraphics[width=0.95\textheight,height=0.92\textwidth,keepaspectratio]{../diagrams/v2/16activity-diagram-uc2.png}
  \caption{Diagrama de actividades — UC2}
\end{sidewaysfigure}

\subsubsection{UC3}
\begin{figure}[H]
  \centering
  \includegraphics[width=\textwidth,height=0.9\textheight,keepaspectratio]{../diagrams/v2/16activity-diagram-uc3.png}
  \caption{Diagrama de actividades — UC3}
\end{figure}

\subsubsection{UC4}
\begin{sidewaysfigure}
  \centering
  \includegraphics[width=0.95\textheight,height=0.92\textwidth,keepaspectratio]{../diagrams/v2/16activity-diagram-uc4.png}
  \caption{Diagrama de actividades — UC4}
\end{sidewaysfigure}

\subsubsection{UC5}
\begin{sidewaysfigure}
  \centering
  \includegraphics[width=0.95\textheight,height=0.92\textwidth,keepaspectratio]{../diagrams/v2/16activity-diagram-uc5.png}
  \caption{Diagrama de actividades — UC5}
\end{sidewaysfigure}

\subsection{Secuencia}
\subsubsection{UC1}
\begin{figure}[H]
  \centering
  \includegraphics[width=0.7\textwidth,height=0.5\textheight,keepaspectratio]{../diagrams/v2/13sequence-diagram-uc1.png}
  \caption{Diagrama de secuencia — UC1}
\end{figure}

\subsubsection{UC2}
\begin{sidewaysfigure}
  \centering
  \includegraphics[width=0.95\textheight,height=0.92\textwidth,keepaspectratio]{../diagrams/v2/13sequence-diagram-uc2.png}
  \caption{Diagrama de secuencia — UC2}
\end{sidewaysfigure}

\subsubsection{UC3}
\begin{sidewaysfigure}
  \centering
  \includegraphics[width=0.95\textheight,height=0.92\textwidth,keepaspectratio]{../diagrams/v2/13sequence-diagram-uc3.png}
  \caption{Diagrama de secuencia — UC3}
\end{sidewaysfigure}

\subsubsection{UC4}
\begin{sidewaysfigure}
  \centering
  \includegraphics[width=0.95\textheight,height=0.92\textwidth,keepaspectratio]{../diagrams/v2/13sequence-diagram-uc4.png}
  \caption{Diagrama de secuencia — UC4}
\end{sidewaysfigure}

\subsubsection{UC5}
\begin{sidewaysfigure}
  \centering
  \includegraphics[width=0.95\textheight,height=0.92\textwidth,keepaspectratio]{../diagrams/v2/13sequence-diagram-uc5.png}
  \caption{Diagrama de secuencia — UC5}
\end{sidewaysfigure}

\newpage
\section{Tolerancia a fallos}
El sistema asume que cualquier nodo del pipeline puede caerse en
cualquier momento —de hecho, durante las pruebas se inyectan caídas con
un \textit{Chaos Monkey}— y debe recuperarse sin perder ni duplicar
resultados. El modelo de fallas considerado es la caída abrupta de un
proceso; se asume que el \textit{broker} RabbitMQ y sus colas durables
sobreviven, y que el volumen de estado de cada nodo persiste a través de
los reinicios.

La garantía de correctitud se construye combinando dos mecanismos.
RabbitMQ entrega \textbf{al menos una vez}: todo mensaje no confirmado
(\textit{ack}) se reenvía tras una caída. Sobre esa base, la
\textbf{deduplicación} descarta los reenvíos ya procesados, de modo que el
efecto observable es un procesamiento \textbf{efectivamente único}. La
detección y el reinicio de los nodos caídos quedan a cargo de un
\textbf{Supervisor} externo, mientras que la consistencia del estado tras el
reinicio la garantiza el \textbf{checkpointing}.

\subsection{Arquitectura de tolerancia a fallos}
Tres piezas sostienen la tolerancia a fallos, ortogonales al pipeline de
procesamiento:

\begin{itemize}
  \item El \textbf{Supervisor}, un proceso externo que recibe los
    \textit{heartbeats} de todos los nodos, detecta los que dejan de
    responder y los revive.
  \item El \textbf{checkpoint} local de cada nodo con estado, que persiste
    en un volumen propio su estado de negocio, su tabla de deduplicación
    y sus contadores de secuencia de salida.
  \item El \textbf{Chaos Monkey}, un proceso de prueba que mata
    contenedores al azar para validar la recuperación.
\end{itemize}

\begin{figure}[H]
  \centering
  \includegraphics[width=\textwidth]{../diagrams/v2/17ft-architecture-diagram.png}\\[1cm]
  \caption{Arquitectura de tolerancia a fallos}
\end{figure}

\subsection{Detección y recuperación de caídas}
Cada nodo levanta, en un \textit{thread} de fondo, un cliente de
\textit{heartbeat} que se registra contra el Supervisor por TCP y le envía
un latido cada segundo. El Supervisor mantiene un registro con el último
latido de cada nodo y, si pasan más de 3 segundos sin recibir uno, marca
al nodo como caído. Un \textit{reviver} recorre periódicamente (cada 2
segundos) los nodos caídos y los reinicia con \texttt{docker start},
reutilizando su volumen de estado; un \textit{cooldown} por nodo evita
reintentos en ráfaga sobre un nodo que tarda en volver. La detección es
puramente por \textit{heartbeat} (reloj monótono, sin sondear Docker ni la
red); sólo la recuperación toca Docker.

\begin{figure}[H]
  \centering
  \includegraphics[width=0.85\textwidth]{../diagrams/v2/18ft-heartbeat-sequence.png}\\[1cm]
  \caption{Detección por heartbeat y revival}
\end{figure}

La línea de tiempo de recuperación, con la configuración de referencia,
es del orden de los 5–6 segundos: tras la caída, en a lo sumo 1 segundo
se pierde el primer latido, a los 3 segundos vence el \textit{timeout} y el
nodo queda marcado como caído, y en la siguiente pasada del \textit{reviver}
(a lo sumo 2 segundos después) se emite el \texttt{docker start}. Al
reiniciarse, el nodo restaura su \textit{checkpoint}, se reconecta a
RabbitMQ y reanuda el consumo.

El cliente de \textit{heartbeat} corre como \textit{daemon} que nunca bloquea
el trabajo del nodo y se reconecta solo si el Supervisor está caído o se
reinicia. El propio Supervisor se despliega con \texttt{restart:
on-failure}, de modo que Docker Compose lo levanta de nuevo si cae;
durante su caída los nodos siguen procesando con normalidad —un nodo que
caiga en esa ventana simplemente no se revive hasta que el Supervisor
vuelve—. No hay elección de líder: la liveness del Supervisor está
desacoplada de la correctitud del pipeline.

\subsection{Fin de stream y barrera de anillo}
El fin de un stream se señaliza con un mensaje especial de \textit{EOF}.
Cuando el Gateway termina de repartir los datos de un cliente, emite un
EOF por cada \textit{shard} con la cantidad exacta de mensajes que envió a
ese shard, de modo que cada instancia downstream sabe cuántos registros
debe haber recibido antes de darse por completa.

El problema aparece cuando una etapa con $N$ instancias en paralelo debe
emitir un único EOF hacia la etapa siguiente: ninguna instancia, por sí
sola, sabe si las demás ya terminaron. Para resolverlo, las instancias de
cada etapa se organizan en un \textbf{anillo} y ejecutan una barrera
distribuida. Cuando una instancia alcanza su conteo esperado, vacía su
estado acumulado hacia downstream y registra cuántos mensajes emitió por
shard. Un \textit{token} de barrera circula por el anillo acumulando los
conteos de cada instancia; cuando vuelve al líder con los aportes de las
$N$, éste emite hacia la etapa siguiente un EOF por shard con el total
agregado del clúster. Así, la etapa downstream recibe exactamente la
cuenta total que debe esperar.

\begin{figure}[H]
  \centering
  \includegraphics[width=0.9\textwidth]{../diagrams/v2/19ft-eof-barrier-diagram.png}\\[1cm]
  \caption{Propagación de EOF mediante la barrera de anillo}
\end{figure}

Este mecanismo es también el que indica a las etapas con estado
(\texttt{aggregate}, \texttt{merge}) cuándo materializar su resultado: una
instancia sólo vacía sus agregados parciales al completarse su entrada.
El \texttt{merge} es un caso particular porque tiene dos entradas (la chica,
por \textit{broadcast}, y la grande, por \textit{sharding}): sólo se da por
completo cuando llegaron los EOF de ambas, usando la suma de los dos
conteos. El Join, en cambio, es de instancia única por caso de uso: no
participa del anillo, sino que espera el EOF agregado y, al alcanzar la
cuenta, materializa y emite las respuestas del cliente.

\subsection{Checkpointing y deduplicación}
Cada nodo con estado persiste periódicamente un \textit{checkpoint} que
incluye su estado de negocio, la tabla de deduplicación y los contadores
de secuencia de salida, todo en un único bloque escrito con
\texttt{msgpack}. La escritura es atómica: se vuelca a un archivo temporal y
se renombra sobre el definitivo, de modo que una caída a mitad de
escritura deja intacto el \textit{checkpoint} anterior.

La pieza clave es el \textbf{orden entre el checkpoint y los acks}. Los
mensajes se procesan en lote (1000 por \textit{checkpoint} en la
configuración de referencia) y sus \textit{acks} se retienen: recién se
confirman a RabbitMQ \textbf{después} de que el \textit{checkpoint} quedó
escrito en disco. Así, todo lo confirmado tiene su estado ya persistido;
si el nodo cae con \textit{acks} pendientes, RabbitMQ reenvía esos mensajes
y se reprocesan sobre el estado restaurado.

\begin{figure}[H]
  \centering
  \includegraphics[width=0.85\textwidth]{../diagrams/v2/20ft-checkpoint-diagram.png}\\[1cm]
  \caption{Checkpoint con acks retenidos y deduplicación}
\end{figure}

La \textbf{deduplicación} es la que evita que ese reproceso duplique
efectos. Cada mensaje de datos se \textit{estampa} en origen con un
identificador de productor estable —derivado de la identidad del nodo y
su ruta de salida, idéntico entre reinicios— y un número de secuencia
creciente. El receptor mantiene, por productor, el mayor número de
secuencia visto; un mensaje cuya secuencia no lo supera es un reenvío y
se descarta. Como bajo ruteo por afinidad un ``productor'' es un par
\texttt{(nodo, ruta)} y no un mensaje, la tabla queda acotada por la
topología. Los contadores de secuencia de salida también se restauran del
\textit{checkpoint}, de modo que un mensaje reemitido tras una caída lleva la
misma secuencia que su original y el deduplicador downstream lo reconoce.
El estado volcado a disco (por ejemplo, los archivos de \textit{spill} del
conteo de caminos de UC4) se trunca, al restaurar, al último \textit{offset}
confirmado, descartando lo que se hubiera escrito después del último
\textit{checkpoint}.

\subsection{Apagado graceful}
Ante una señal \texttt{SIGTERM} (o \texttt{SIGINT}), cada nodo corta el
consumo, cierra su conexión con RabbitMQ y libera sus publicadores. Los
\textit{acks} del lote en curso no se confirman: al cerrarse la conexión,
RabbitMQ reencola esos mensajes y el próximo arranque los reprocesa
sobre el último \textit{checkpoint}. El apagado es, por tanto, equivalente
en correctitud a una caída: el sistema cumple el \textit{graceful quit}
exigido por el enunciado sin necesidad de vaciar trabajo en vuelo.

\subsection{Puntos de falla}
La combinación de entrega al-menos-una-vez, deduplicación, \textit{acks}
retenidos y barrera idempotente cubre los distintos momentos en que un
nodo puede caer. A continuación se enumeran los puntos protegidos y cómo
se preserva la correctitud en cada uno.

\subsubsection{Caída a mitad de lote}
Un nodo procesa parte de un lote (aplicando estado y registrando
deduplicación en memoria) y cae antes de \textit{checkpointear}. Como esos
mensajes no se confirmaron, RabbitMQ los reenvía; el estado y la
deduplicación se restauran del último \textit{checkpoint} y los mensajes se
reprocesan limpiamente sobre él.

\begin{figure}[H]
  \centering
  \includegraphics[width=0.85\textwidth]{../diagrams/v2/21ft-failure-midbatch.png}\\[1cm]
  \caption{Caída a mitad de lote}
\end{figure}

\subsubsection{Caída tras el checkpoint y antes de los acks}
El nodo escribe el \textit{checkpoint} del lote y cae antes de confirmar los
\textit{acks}. RabbitMQ reenvía el lote completo, pero al restaurar, la tabla
de deduplicación ya contiene esas secuencias: cada mensaje reenviado se
detecta como duplicado, se confirma y se descarta sin volver a aplicarse.
El estado queda correcto, tomado del \textit{checkpoint}.

\begin{figure}[H]
  \centering
  \includegraphics[width=0.85\textwidth]{../diagrams/v2/22ft-failure-preack.png}\\[1cm]
  \caption{Caída tras el checkpoint, antes de los acks}
\end{figure}

\subsubsection{Caída de un nodo con estado acumulado}
Un \texttt{aggregate} o \texttt{merge} cae con resultados parciales
acumulados a lo largo de varios \textit{checkpoints}. Se restaura el estado
del último \textit{checkpoint} y, si ya había emitido y enviado el
\textit{token} de barrera, la barrera es idempotente: re-recibir el token
sólo vuelve a registrar los mismos conteos, sin doble emisión. Los
archivos de \textit{spill} se truncan a su \textit{offset} confirmado.

\begin{figure}[H]
  \centering
  \includegraphics[width=0.9\textwidth]{../diagrams/v2/23ft-failure-stateful.png}\\[1cm]
  \caption{Caída de un nodo con estado acumulado}
\end{figure}

\subsubsection{Caída del Gateway}
Es el límite de diseño conocido. El Gateway asigna un \texttt{client\_id}
nuevo en cada conexión, por lo que si cae mientras un cliente transmitía,
los datos parciales ya publicados quedan huérfanos: ningún EOF llegará
para ese \texttt{client\_id} y el Join nunca finaliza esa sesión. El cliente
debe reconectarse y reenviar el trabajo desde cero. Es el único escenario
en el que la correctitud depende del reintento del cliente.

\begin{figure}[H]
  \centering
  \includegraphics[width=0.9\textwidth]{../diagrams/v2/24ft-failure-gateway.png}\\[1cm]
  \caption{Caída del Gateway}
\end{figure}

\subsubsection{Caída del Supervisor}
Docker Compose lo reinicia (\texttt{restart: on-failure}). Los nodos siguen
procesando sin interrupción; si alguno cae durante la ventana de caída
del Supervisor, no se revive hasta que éste vuelve y los nodos se
re-registran. No hay falla en cascada: la liveness del Supervisor está
desacoplada de la correctitud del pipeline.

\begin{figure}[H]
  \centering
  \includegraphics[width=0.85\textwidth]{../diagrams/v2/25ft-failure-supervisor.png}\\[1cm]
  \caption{Caída del Supervisor}
\end{figure}

\newpage
\section{Desarrollo}
\subsection{Tareas a realizar}
Para el desarrollo del trabajo, se decidió que cada uno de los
integrantes realice la implementación de uno o más casos de uso de
punta a punta; con el objetivo de que todos podamos tener un
seguimiento del total funcionamiento del sistema, pero sin la
necesidad de entrar de lleno en los detalles específicos de cada caso de uso.

\subsection{Asignación}
Aquellas tareas de las que dependan todas los casos, serán
desarrolladas en una primera implementación por todos los miembros
del equipo, y serán modificadas como se requiera por quien lo
necesite a lo largo de la continuación del trabajo.
El primer caso de uso será desarrollado en gran medida con la
dinámica de *ping-pong pair-programming*; con el objetivo de usar el
caso más simple para la interiorización de la implementación del
trabajo para todos los miembros del equipo en paralelo.
El resto de ítems quedan divididos de la siguiente manera:

\begin{itemize}
  \item UC1: Todos
  \item UC2: Lorenzo Minervino
  \item UC3: Valsagna Federico.
  \item UC4: Ordoñez Alejo.
  \item UC5: Lorenzo Minervino.
\end{itemize}

\end{document}
